% lejepa_deeper_look.tex — appendix "A closer look at LeJEPA" (draft 2026-09-17, station = the embedding per Berker; RAW,
% not jointly agreed). Replaces the stub \paragraph{Each method's own regularizer at the backbone.} in iclr2027_conference.tex.
% Runs: vanilla = in100.lejepa.s0 (E17 control re-extraction e17c.ext); + SIGReg at h = in100.lejepa.s0.hpull_sigreg.ext
% (E17 arm, h_lamb .003763 = 10 % of the projector-side SIGReg's weighted backbone pull on the first batch).
% Station = student.z.embed = LeJEPA's encoder output (ViT(num_classes=512): the 512-d Linear is inside their encoder;
% docs/methods/lejepa.md l.105; D-036 declared h for the lane; eval_features returns it).
% Sources: pos/neg cosine, moment KL (full covariance, gauss_kl_full: D-040), RankMe: results/diag/e17_depth_metrics.csv
% (labels lejepa_ctrl / lejepa_sigreg, space student.z.embed); Epps–Pulley: results/battery/<run>.csv (train500.v1L, raw|full,
% n = 50k both); per-coordinate KL: E17 card §First scoring pass ("diag-KL@h .068 (.956)"; archive/e17_score.py::diag_kl =
% 0.5·(mean μ_i² + mean(σ_i² − 1 − log σ_i²)) over coordinates); probes: results/probes/<run>.csv (linear_raw_v2, knn_v1_k200).
% Our term (text): in100.lejepa.s0.e20f (E20; e20_zoo_depth_metrics.csv, battery e20f.ext.csv, probes e20f.ep100.extL.csv).
\section{A closer look at LeJEPA}
\label{app:lejepa}

LeJEPA is the closest method to ours in spirit. It argues that the ideal representation is an isotropic Gaussian and enforces this with SIGReg, a sketched test of Gaussianity along random one-dimensional projections, applied after the projector together with the invariance term. The representation it keeps, the encoder output, is far from that ideal (Figure~\ref{fig:hz_geometry}). We therefore move the loss to where it matters and ask whether SIGReg itself works as a backbone regularizer. Starting from vanilla LeJEPA (SIGReg and invariance after the projector, both unchanged), we add a copy of SIGReg on the encoder output, the 512-dimensional representation that feeds the projector, with a small weight ($\lambda=0.0038$, set so that the added term's gradient on the backbone is 10\% of the projector-side SIGReg's at initialization), and train under the ImageNet-100 setup of Section~\ref{sec:exp-in100}. If SIGReg were the right backbone regularizer, it should push the representation toward LeJEPA's own ideal, an isotropic Gaussian, and downstream performance should follow.

Table~\ref{tab:lejepa-sigreg} reports the readouts at the encoder output, where the term acts. By its own measures SIGReg does make the representation more Gaussian: the mean offset that made unrelated images similar disappears (negative-pair cosine $0.58\to0.01$), every coordinate becomes zero-mean and unit-variance (per-coordinate KL $0.96\to0.07$), and the Epps--Pulley statistic falls. But the joint distribution does not follow: the moment KL that accounts for the full covariance stays high ($3.89\to3.03$) because the spectrum does not broaden (RankMe$/d$ $0.07\to0.09$), the views of an image move apart (positive-pair cosine $0.96\to0.73$), and both probes fall (linear $-1.4$, kNN $-4.3$). Two conclusions follow. First, more Gaussianity at the representation does not improve performance by itself. Second, SIGReg does not work as an alternative backbone regularizer: under the same setup and at the same tap, our regularizer raises RankMe$/d$ to $0.45$, lowers the full-covariance KL to $1.73$, and raises the probes by $+6.6$ (linear) and $+10.5$ (kNN). The difference lies in what the two terms can see. SIGReg matches one-dimensional projections to a Gaussian, and these look Gaussian as soon as the mean is removed and the coordinates are standardized, whatever the number of directions that carry variance. The log-determinant in $\mathcal R_{\mathrm{ctr}}$ instead penalizes every direction of the covariance whose variance vanishes.

\begin{table}[h]
\centering
\small
\begin{tabular}{lrr}
\toprule
Readout at the encoder output ($h$) & LeJEPA & LeJEPA + SIGReg at $h$ \\
\midrule
Positive-pair cosine & 0.96 & 0.73 \\
Negative-pair cosine & 0.58 & 0.01 \\
Per-coordinate KL to $\mathcal N(0,1)$ & 0.96 & 0.07 \\
Moment KL to $\mathcal N(0,I)$ (full covariance) & 3.89 & 3.03 \\
Epps--Pulley & 652 & 504 \\
RankMe$/d$ & 0.07 & 0.09 \\
Linear probe (\%) & 60.4 & 58.9 \\
kNN, $k=200$ (\%) & 52.4 & 48.1 \\
\bottomrule
\end{tabular}
\caption{Vanilla LeJEPA against LeJEPA with SIGReg added at its encoder output (ImageNet-100, ViT-S/16, 100 epochs; projector losses unchanged). All readouts at the 512-dimensional encoder output that LeJEPA keeps as its representation: pair cosines on eight stored views of 10{,}000 training images; the per-coordinate KL uses each coordinate's mean and variance, the moment KL the full covariance (its log-determinant sees vanishing directions); Epps--Pulley on the clean training subset (same sample size for both models); RankMe over the feature dimension; the linear and kNN probes of Section~\ref{sec:exp-setup}.}
\label{tab:lejepa-sigreg}
\end{table}
% Rank recheck (Berker 2026-09-17 "25-30 effective dimensions and still 60 %?"; scratch/lejepa_rank_recheck.py, job 66273881,
% outputs/e-rankcheck_66273881.out; raw features of features/in100.lejepa.s0.ext, converged L-BFGS probes on standardized PCA
% coordinates): LeJEPA embedding — 98.4 % of the variance in the top 32 directions, 45 eigenvalues above 1e-3 of the largest,
% RankMe 35 uncentered / 43 centered; probes on the top 8/16/32/64/128 directions = 39.7/50.8/57.6/60.7/62.5 (full feature, paper
% reader: 60.4). LeJEPA CLS — top 32 = 91.8 % of the variance; tail beyond 64 (2.9 % of the variance) alone gives 51.2. Ours at h,
% embedding: RankMe 229, top 32 = 26.9 % of the variance, top-64/128 probes 60.5/64.6, tail beyond 64 (54.6 % of the variance) 52.8.
% So the low rank is real and 60 % fits in ~32-64 directions; entropy ranks count variance, not informative directions.
% RAW (not jointly agreed). §2's figure now reads LeJEPA and VISReg at their embedding too (H_STATION rule in paper_exhibits.py,
% 2026-09-17), so the RankMe/d 0.07 here matches it; the §2 TEXT still quotes the old CLS values (0.23 / 0.28). Same story at the CLS: moment KL 3.74 -> 3.31,
% RankMe/d 0.23 -> 0.24, positive cosine 0.97 -> 0.73, negative 0.77 -> 0.01, Epps–Pulley 522 -> 275, linear 65.5 -> 62.2,
% kNN 51.4 -> 48.6 (per-coordinate KL at the CLS not available: the E17 stores were purged; both checkpoints still exist).
